\section{Introduction}
\label{sec:intro}
Living donor kidney transplantation (LDKT) remains the treatment of choice for end-stage renal disease (ESRD), offering better graft function and longer allograft survival than deceased-donor transplantation.~\cite{srtr_2023} Yet the gap between organ need and organ supply persists: more than 4,000 candidates died on the US kidney waiting list in 2024, and a comparable number were removed for transplant-disqualifying deterioration, while the annual number of LDKTs has remained essentially flat for over a decade.~\cite{cooper_100-year_2026} Expanding living donation therefore rests on two complementary levers: (1) recruiting more donors and (2) better understanding, quantifying, and managing the risk that donation poses to the donor.

The safety literature on living donation is large, but has historically emphasized relative rather than absolute, or life-expectancy-scaled, risk. Perioperative mortality is low (3.1 per 10,000 donors) but not uniform, rising sharply with male sex, Black race, and pre-existing hypertension.~\cite{segev_perioperative_2010} Longer-term follow-up using healthy, contraindication-screened non-donor comparators, rather than the general population, established that living donors face a several-fold higher risk of subsequently developing ESRD themselves, even though the absolute magnitude of that risk remains small.~\cite{muzaale_risk_2014, massie_quantifying_2017} Whether this translates into excess all-cause mortality is unresolved: a meta-analysis of 52 studies and over 118,000 donors found no elevation in all-cause mortality relative to matched non-donors,~\cite{okeeffe_mid-_2018} whereas a Norwegian registry study with unusually long follow-up and a cohort drawn predominantly from first-degree relatives of recipients reported a modestly elevated hazard ratio of 1.30.~\cite{mjoeen_long-term_2014} This tension between the largest available meta-analysis and the longest single-cohort follow-up is a central source of uncertainty for any model of donor life expectancy.

Post-donation ESRD risk is not distributed evenly across the donor population. Black donors experience roughly three times the 15-year cumulative incidence of ESRD observed in White donors, a disparity that persists after matching to equally healthy non-donor comparators of the same race.~\cite{muzaale_risk_2014} Part of this excess risk has a defined biological correlate: two-copy \textit{APOL1} renal-risk-variant genotypes, disproportionately carried by individuals of recent African ancestry, are among the strongest known risk factors for progression to CKD stage~3 or higher and have been directly implicated in the long-term kidney health of Black living donors.~\cite{doshi_apol1_2018} Donation itself is also not evenly distributed: women have consistently made up a majority of living donors, a pattern attributable in part to socioeconomic and caregiving asymmetries, financial barriers to donation, and disparities in access to unpaid leave that fall more heavily on lower-income men.~\cite{matas_gender_2018} Together, these findings motivate stratifying any life-expectancy analysis of donation by age, race, and sex rather than reporting a single population-average effect.

Policy has responded to donor risk primarily by protecting donors' future access to a transplant should they themselves develop ESRD. The Kidney Allocation System implemented in 2014 (KAS250) grants prior living donors priority access to the deceased-donor waiting list, and early evaluation confirmed that this priority meaningfully shortened prior donors' time to listing.~\cite{wainright_impact_2017} Independently, a national registry analysis of transplant candidates who had previously donated an organ found that such priority access translated into substantially shorter waiting times and higher-quality kidneys compared with matched non-donor candidates, confirming that the policy achieves its intended effect at the level of the individual recipient.~\cite{reese_kidney_2015} The National Kidney Registry has extended this priority concept beyond the donor's own future risk through its voucher program, under which a donor may designate family members, or in some cases any future intended recipient, to receive expedited, priority access to a living-donor kidney should they ever need one. Two large Monte Carlo simulation studies, spanning 50-year and, subsequently, 100-year horizons, have shown that the voucher program is programmatically sustainable: under a wide range of donor-growth scenarios, the supply of future voucher-eligible donors is projected to exceed voucher redemptions by a comfortable margin.~\cite{cooper_ensuring_2021, cooper_100-year_2026}

These two literatures, quantifying the donor's own excess ESRD and mortality risk, and demonstrating the programmatic feasibility of priority-access mechanisms such as KAS250 and the NKR voucher, have so far remained largely separate. The voucher simulations answer a supply-side question (will there be enough donors to honor future redemptions?) but do not quantify the demand-side question of how much life expectancy a voucher, or PLD priority access more generally, is actually worth to the person who holds it. Similarly, the ESRD-risk literature reports relative risks and cumulative incidences but stops short of translating these into an integrated, individual-level life-expectancy trade-off that simultaneously accounts for the donor's elevated ESRD risk, the mortality and quality-of-life consequences of dialysis and transplant waitlisting, and demographic heterogeneity. This paper addresses that gap using a Markov microsimulation model that estimates the change in life expectancy attributable to kidney donation itself, stratified by age at donation, race, and sex, and separately quantifies the life-expectancy value of priority waitlist access, both unconditionally at the population level and conditional on the recipient having actually developed ESRD.